

\section{Différents modèles étudiés}

\begin{frame}
	\frametitle{Différents modèles}
	...
\end{frame}

%=====================================================================================================

\subsection{Modèle 1}

	\begin{frame}{Méthode analytique}
		
		
		\begin{block}{Avantages}
			\begin{enumerate}
				\item Applicable à tous les avions.
				\item Les forces sont calculées en fonction de la vitesse à chaque instant.
				\item Peut être utilisé pour tout type d'avion.
			\end{enumerate}
		\end{block}
		
		\begin{block}{Hypothèses}
			\begin{enumerate}
				\item La vitesse de décollage $V_{LOF}$ est connue à l'avance.
				\item Le moteur est supposé avoir développé toute sa poussée avant le départ.
				\item La masse est supposée constante.
			\end{enumerate}
		\end{block}
		
	\end{frame}
	
	%----------------------------------------------------------------------------------------------------------------
	
	\begin{frame}
		
		\frametitle{Méthode}
		
		En utilisant le PFD de l'équation \ref{PFD} et la relation de dérivation totale :
		
		\begin{equation}
			\frac{ dV}{dt}=\frac{ds}{dt}\frac{dV}{ds}=(V-V_{hw})\frac{dV}{ds}
		\end{equation}
		
		On obtient alors :
		
		\begin{equation}
			\frac{ds}{dV}=\frac{W(V-V_{hw})}{g(T-D-F_r)}
		\end{equation}
		
		En intégrant entre deux vitesses successives :
		
		\begin{align}
			S_{i+1}-S_i=
			\frac{1}{g}
			\int_{V_i}^{V_{i+1}}
			\frac{V-V_{hw}}
			{\dfrac{T-D-F_r}{W}}
			\, dV
		\end{align}
		
	\end{frame}
	
	%----------------------------------------------------------------------------------------------------------------
	
	\begin{frame}{Modélisation des forces}
		
		Dans ce modèle, les forces sont de la forme :
		
		\begin{align}
			D &= \frac{1}{2}\rho V^2 S_w C_D \\
			F_r &= \mu_r W - \frac{1}{2}\mu_r \rho S_w C_L V^2 \\
			T &= (T_0)_i + (T')_i V + (T'')_i V^2 
		\end{align}
		Avec $ (T_0)_i , (T')_i , (T'')_i $ qui sont des constantes expérimentales
		
	\end{frame}
	
	%=====================================================================================================
	
	%=====================================================================================================
\subsection{Modèle de forces constantes}

	\begin{frame}
		\frametitle{Modèle de forces constantes}
		
		\begin{block}{Avantages}
			\begin{enumerate}
				\item Applicable à tous les avions si les forces peuvent être approximées.
				\item Applicable aux avions à hélice et aux turboréacteurs.
			\end{enumerate}
		\end{block}
		
		\begin{block}{Hypothèses}
			\begin{enumerate}
				\item Les forces sont constantes pendant le roulage sur la piste.
				\item Les forces sont estimées à la vitesse $\dfrac{V_R}{\sqrt{2}}$.
				\item La masse est supposée constante.
			\end{enumerate}
		\end{block}
		
	\end{frame}
	
	%----------------------------------------------------------------------------------------------------------------
	
	\begin{frame}
		
		\frametitle{Modèle de forces constantes: Méthode}
		
		En utilisant le PFD de l'équation \ref{PFD} :
		
		\begin{equation*}
			S_G=
			\frac{V_R^2 W}
			{2g[T-D-\mu(W-L)]}
			\quad
			\text{évaluée en }
			\left(\frac{V_R}{\sqrt{2}}\right) 
		\end{equation*}
		
		avec $ V_{R}$ la vitesse de rotation
		
	\end{frame}
	
	%----------------------------------------------------------------------------------------------------------------
	
	\begin{frame}{Modélisation des forces}
		
		Dans ce modèle, seule la force de poussée est différente des autres forces.
		
		\textbf{Moteur à piston}
		\[
		T=\frac{n_p(550P_{BHP})}{V_R/\sqrt{2}}
		\]
		
		\textbf{Moteur à turbopropulseur}
		\[
		T=T\left(\frac{V_R}{\sqrt{2}}\right)
		\]
		
	\end{frame}
	
	%----------------------------------------------------------------------------------------------------------------
	
	\begin{frame}{Cas particulier de la méthode}
		
		Un cas particulier de cette méthode consiste à déterminer la distance de décollage pour un avion à train tricycle.
		
		\begin{block}{Avantages}
			\begin{enumerate}
				\item Applicable uniquement aux avions à train tricycle.
				\item Permet une meilleure étude des paramètres de l'équation.
			\end{enumerate}
		\end{block}
		
		Dans ce cas :
		
		\begin{equation*}
			S_G=
			\frac{V_R^2 W}
			{n_p\,550\,P_H
				+16.09\rho V_R^3 S(\mu C_{LTO}-C_{DTO})
				-2g\mu W V_R}
		\end{equation*}
		
	\end{frame}
	
	
	%=====================================================================================================
	%=====================================================================================================
	
	
		\subsection{Modèle 3}
		
			\begin{frame}
				\frametitle{Méthode numérique}
				
				\begin{block}{Avantages}
					\begin{enumerate}
						\item Applicable à tous les avions.
						\item Prend en compte la distance de freinage.
						\item Utile pour les avions à train classique (taildragger) qui passent de trois à deux roues.
					\end{enumerate}
				\end{block}
				
				\begin{block}{Hypothèses}
					\begin{enumerate}
						\item La vitesse de décollage est supposée égale à $V_{LOF} = 1.1 V_{s1}$.
						\item Le moteur est supposé avoir développé toute sa poussée avant le départ.
						\item La masse est supposée constante.
						\item La piste est supposée horizontale et les conditions atmosphériques sont standard.
					\end{enumerate}
				\end{block}
				
			\end{frame}
			
			%----------------------------------------------------------------------------------------------------------------
			
			\begin{frame}
				
				\frametitle{Méthode}
				
				\textbf{Étape 2 : Calcul de l'accélération}
				
				\[
				a_i=\frac{g}{W}\left[T_i-D_i-\mu(W-L_i)\right]
				\]
				
				\textbf{Étape 3 : Mise à jour de la vitesse et de la distance}
				
				\[
				V_i=V_{i-1}+a_i\Delta t_i
				\]
				
				\[
				S_i=S_{i-1}+V_{i-1}\Delta t_i+\frac{1}{2}a_i\Delta t_i^2
				\]
				
				avec
				
				\[
				\Delta t_i=t_i-t_{i-1}
				\]
				
				\textbf{Étape 4 : Condition de décollage}
				
				Le calcul est répété jusqu'à :
				
				\[
				V_i=V_{LOF}
				\]
				
				La distance obtenue est :
				
				\[
				S_i=S_G
				\]
				
				représentant la distance totale de roulement au sol.
				
			\end{frame}
			
			%----------------------------------------------------------------------------------------------------------------
			
			\begin{frame}{Modélisation des forces}
				
				À chaque pas de temps $i$, on calcule :
				
				\begin{itemize}
					\item la poussée : $T_i$ ;
					\item la traînée : $D_i$ ;
					\item la portance : $L_i$ ;
					\item la densité de l'air : $\rho$ ;
					\item la vitesse au pas précédent : $V_{i-1}$.
				\end{itemize}
				
			\end{frame}
			
			
			
		

	
	
	%=====================================================================================================
\subsection{Modèle 4}
	
	\begin{frame}
		\frametitle{Modèle 4}
		
		\begin{block}{Avantages}
			\begin{enumerate}
				\item Prise en compte de la dégradation de la poussée ($k$).
				\item Ne nécessite ni boucle de calcul ni pas de temps $\Delta t$.
			\end{enumerate}
		\end{block}
		
		\begin{block}{Hypothèses}
			\begin{enumerate}
				\item Les coefficients aérodynamiques ($C_L$ et $C_D$) sont moyennés.
				\item La propulsion nette dépend directement de la vitesse aérodynamique instantanée ($V_\infty$).
				\item La force de propulsion est maximale à l'arrêt.
				\item La constante $k$, dépendant du moteur, est prise en compte.
			\end{enumerate}
		\end{block}
		
	\end{frame}
	
	%----------------------------------------------------------------------------------------------------------------
	
	\begin{frame}
		
		\frametitle{Méthode}
		
		Reprenons l'équation de l'accélération \ref{PFD} projetée selon l'axe horizontal :
		
		\begin{equation}
			\frac{dv}{dt}
			=
			\frac{\left(T-\mu_r(W-L)-D\right)g}{W}
		\end{equation}
		
	\end{frame}
	
	%----------------------------------------------------------------------------------------------------------------
	
	\begin{frame}{Modélisation des forces}
		
		Dans ce modèle, les forces sont de la forme :
		
		\begin{align}
			D &= \frac{1}{2}\rho V^2 S_w C_D \\
			F_r &= \mu_r W - \frac{1}{2}\mu_r \rho S_w C_L V^2 \\
			T &= T_0-kV^2
		\end{align}
		
	\end{frame}
	
	%----------------------------------------------------------------------------------------------------------------
	
	\begin{frame}
		
		\frametitle{Développement}
		
		En utilisant la formule de la dérivation totale, on obtient :
		
		\begin{equation}
			V_{\infty}\frac{dV_{\infty}}{ds}
			=
			A-BV_{\infty}^2
		\end{equation}
		
		où les constantes scalaires $A$ et $B$ s'expriment de la façon suivante :
		
		\begin{equation}
			A=
			\left(
			\frac{T_0}{W}-\mu_r
			\right)g
			\qquad
			\text{et}
			\qquad
			B=
			\frac{g}{W}
			\left[
			\frac{1}{2}\rho_{\infty}S(C_D-\mu_rC_L)+k
			\right]
		\end{equation}
		
		Séparons les variables en regroupant tout ce qui dépend de $V_{\infty}$ d'un côté et $ds$ de l'autre :
		
		\begin{equation}
			ds=
			\frac{V_{\infty}\,dV_{\infty}}
			{A-BV_{\infty}^2}
		\end{equation}
		
	\end{frame}
	
	%=====================================================================================================
	\subsection{Modèle 5}
	
		\begin{frame}
			\frametitle{Modèle 5}
			Développement d'une nouvelle méthode proposée par l'équipe...
		\end{frame}